\documentclass[10pt,a4paper]{amsart}
\usepackage[margin=19mm]{geometry}
\usepackage{amsmath,amssymb,mathtools}
\usepackage[scr=rsfs,cal=euler]{mathalfa}
\usepackage{xfrac}
\usepackage[unicode,hidelinks,bookmarks=false]{hyperref}
\newcommand{\ie}{i.e.\ }
\newcommand{\cf}{cf.\ }
\newcommand{\resp}{respectively\ }
\newcommand{\Subsection}{\subsection}
\providecommand{\nobreakdash}{\nobreak\hbox{-}}
\setlength{\emergencystretch}{2em}
\multlinegap0pt
\usepackage{bm}
\usepackage{bbm}
\usepackage{dsfont}
\usepackage{xspace}
\usepackage{cancel}
\usepackage{mathtools}
\usepackage[shortlabels]{enumitem}
\usepackage{amsthm}
\makeatletter
\@ifundefined{lemma}{\newtheorem{lemma}{Lemma}}{}
\@ifundefined{theorem}{\newtheorem{theorem}{Theorem}}{}
\makeatother
\makeatletter
\expandafter\ifx\csname AMS\endcsname\relax
\newenvironment{AMS}{}{}
\fi
\expandafter\ifx\csname proof\endcsname\relax
\newenvironment{proof}{{\sc Proof:}}{~\hfill QED}
\fi
\makeatother
\title[Surfaces associated with first-order ODEs]{Surfaces associated with first-order ODEs}
\author{Antonio J. Pan-Collantes, Jos\'e A. \'Alvarez-Garc\'ia}
\date{Source: arXiv:2312.04489v4 (CC BY 4.0)}
\begin{document}
\maketitle

\noindent\emph{Source and licence.} Surfaces associated with first-order ODEs. Version arXiv:2312.04489v4 (CC BY 4.0), distributed under the Creative Commons Attribution 4.0 International licence, as recorded in the arXiv licence metadata for this exact version and shown on the abstract page https://arxiv.org/abs/2312.04489v4. Full rights notice: licenses/arxiv-2312.04489-CC-BY-4.0.txt.\par\smallskip

\noindent\textit{Editorial scope.} Counted unit: the Theorem whose source label is \texttt{th:Jacobi} in the article (source0.tex lines 1136--1138, source label \texttt{th:Jacobi}), delivered together with the article's own complete original proof of it (source0.tex lines 1140--1184, one \texttt{proof} environment of 45 lines). No mathematics is written, repaired or re-derived: statement and proof are the article's own text line for line. Required context retained uncounted: ODE (lines 232--234); 1form (lines 258--260); LCdef (lines 1011--1016); prop:intsymop (lines 1063--1069); lem:op (lines 1113--1118); JacobiDef (lines 1130--1132). Every definition, display and label of those lines is retained unchanged. Omitted from this package and not redistributed: 1--231, 235--257, 261--1010, 1017--1062, 1070--1112, 1119--1129, 1133--1135, 1139--1139, 1185--1350. Nothing inside a retained excerpt is omitted or altered: every retained body is delivered exactly as the article's lines are written, so no omission receipt of any kind accompanies this package. Citations: the retained excerpts carry no citation command at all, so no bibliography entry of the article is transcribed and no reference is redistributed. Reference adaptation: none was needed and none was made; every reference inside the delivered unit resolves inside it, so no unresolved reference or citation is produced. Printed slips kept literal: no printed word, symbol, hypothesis or number of the counted statement or of its proof was altered. Layout and class: the article's own class is \texttt{article} with packages including PRIMEarxiv, comment, inputenc, fontenc, url, booktabs, amsfonts, nicefrac, microtype, fancyhdr, graphicx, todonotes, xcolor, tikz-cd; the wrapper is the lane's pinned \texttt{amsart} editorial wrapper at 10pt on A4, which restates the theorem-like environments the retained text needs and adds the article's own macro definitions that the retained text needs (none), with extra wrapper packages (bm, bbm, dsfont, xspace, cancel, mathtools, enumitem). The delivered numbering is the unit's own. Assets: no retained span contains a figure environment, a graphics-inclusion command, a PDF-page inclusion, a table or a TikZ picture; no image file is copied into this package, the package declares no asset dependency and no image file is redistributed with it. Licence: version arXiv:2312.04489v4 of this article is distributed under the Creative Commons Attribution 4.0 International licence, as recorded in the arXiv licence metadata for this exact version and shown on the abstract page https://arxiv.org/abs/2312.04489v4. A full rights notice is delivered at licenses/arxiv-2312.04489-CC-BY-4.0.txt. Characters: every retained excerpt is pure ASCII, so the delivered engine, XeLaTeX with its default Latin Modern text font, reports no missing character.
\begin{equation}\label{ODE}  
    \frac{du}{dx}=\phi(x,u),  
\end{equation}  

\begin{equation}\label{1form}  
    \omega=-\phi(x,u)dx+du.
\end{equation}

\begin{equation}\label{LCdef}
\Theta=\begin{pmatrix}
    0 & -\Delta_{\epsilon} \omega_{\epsilon}^2\\ 
    \Delta_{\epsilon} \omega_{\epsilon}^2 & 0
\end{pmatrix}.    
\end{equation}

\begin{lemma}\label{prop:intsymop}
Let $h\in \mathcal{C}^{\infty}(U)$. Then:
\begin{enumerate}[label=(\alph*)]
    \item \label{it:integrating} $\mathfrak{T}_{\epsilon}(h)=0$ if and only if $e^{\epsilon} h$ is an integrating factor of \eqref{1form}.
    \item If $h$ is non-vanishing in $U$, then $\mathfrak{S}_{\epsilon}(h)=0$ if and only if $e^{\epsilon} h^{-1}$ is an integrating factor of \eqref{1form}.
\end{enumerate} 
\end{lemma}

\begin{lemma}\label{lem:op}
    It is satisfied that
    \begin{equation}
        \mathfrak{T}_{\epsilon}\circ \mathfrak{S}_{\epsilon}=A^2+\mathcal{K}_{\epsilon}.
    \end{equation}
\end{lemma}

\begin{equation}\label{JacobiDef}
    \nabla_{A} \nabla_{A}J+\mathcal{K}_{\epsilon}\left(J-g(J,A)A\right)=0.
\end{equation}

\begin{theorem}\label{th:Jacobi}
The knowledge of a non-trivial Jacobi vector field relative to $A$ on any of the surfaces $\mathcal{S}_{\epsilon}$ leads to the integration of the first-order ODE \eqref{ODE}.
\end{theorem}

\begin{proof}
Let $J$ be a non-trivial Jacobi field relative to $A$ on $\mathcal{S}_{\epsilon}$, for certain $\epsilon\in \mathcal{C}^{\infty}(U)$. We can decompose $J$ into its components with respect to the frame $\{A,e^{-\epsilon}\partial_u\}$ as follows:
$$
J=\sigma A+\delta e^{-\epsilon}\partial_u,
$$
where $\sigma=g_{\epsilon}(J,A)$ and $\delta=g_{\epsilon}(J, e^{-\epsilon}\partial_u)$. 

Consider, first, the case where $\delta\neq 0$. We assume, without loss of generality, that $\delta$ is non-vanishing in $U$ (if necessary, $U$ can be appropriately reduced). Since $J$ is a Jacobi field relative to $A$, equation \eqref{JacobiDef} holds:
\begin{equation}
    \nabla_{A} \nabla_{A}\left( \sigma A+\delta e^{-\epsilon}\partial_u\right)+\mathcal{K}_{\epsilon} \delta e^{-\epsilon}\partial_u=0.
\end{equation}
Now, since $\nabla_{A} A=\nabla_{A}\left(  e^{-\epsilon}\partial_u\right)=0$, as can be deduced from the connection 1-form \eqref{LCdef}, we obtain 
\begin{equation}
    A^2(\sigma) A+A^2(\delta) e^{-\epsilon}\partial_u+ \mathcal{K}_{\epsilon} \delta e^{-\epsilon}\partial_u=0.
\end{equation}
Therefore, $\delta$ must satisfy
\begin{equation}\label{eq:JacDelta}
A^2(\delta)+\mathcal{K}_{\epsilon} \delta=0.
\end{equation}

According to Lemma \ref{lem:op}, we have that
$$
\left(\mathfrak{T}_{\epsilon}\circ \mathfrak{S}_{\epsilon} \right)(\delta)=0,
$$
and we can use Lemma \ref{prop:intsymop} to conclude that
\begin{itemize}
    \item either $\mathfrak{S}_{\epsilon}(\delta)=0$ and then $e^{\epsilon} \delta^{-1}$ is an integrating factor of \eqref{ODE};
    \item or $\mathfrak{S}_{\epsilon}(\delta)\neq 0$ and $\mathfrak{T}_{\epsilon}(\mathfrak{S}_{\epsilon}(\delta))=0$, so $e^{\epsilon} \mathfrak{S}_{\epsilon}(\delta)$ is an integrating factor for \eqref{ODE}.
\end{itemize}
Therefore, equation \eqref{ODE} can integrate it by quadratures.

Suppose now that $\delta=0$, and therefore $\sigma\neq 0$. If we write equation \eqref{JacobiDef} in this particular case we obtain:
\begin{equation}
    \nabla_{A} \nabla_{A}\left( \sigma A\right)=A^2(\sigma)A=0,
\end{equation}
and $A^2(\sigma)=0$. Consequently, if $A(\sigma)$ is not constant, it is a non-trivial first integral of $A$. 

On the contrary, if $A(\sigma)= a\in \mathbb R$, then $A(\sigma-ax)=0$, and $\sigma-ax$ would be a first integral of $A$. Observe that it is, indeed, a non-trivial first integral, because if 
$$
\sigma-ax=b, b\in\mathbb R,
$$
then $J=(ax+b)A$, and it would be a trivial Jacobi field relative to $A$.

In any case, equation \eqref{ODE} is solved, and the result is proven.
\end{proof}

\end{document}
